% ECONOMICS_PROBLEM: {"id": "D73-170", "title": "Effective bureaucratic organizations", "jel_code": "D73", "source_id": "OP-0382"}
% !TeX program = xelatex
\documentclass[11pt,a4paper]{article}
\usepackage[margin=23mm]{geometry}
\usepackage{fontspec}
\setmainfont{Latin Modern Roman}
\usepackage{amsmath,amssymb,mathtools}
\usepackage{xcolor,enumitem,booktabs,longtable,array,xurl,hyperref}
\definecolor{muted}{HTML}{53636C}
\hypersetup{colorlinks=true,urlcolor=blue,linkcolor=blue}
\setlength{\parindent}{0pt}
\setlength{\parskip}{5pt}
\newcommand{\R}{\mathbb R}
\newcommand{\E}{\mathbb E}
\newcommand{\N}{\mathbb N}
\newcommand{\ind}{\mathbf 1}
\newcommand{\Var}{\operatorname{Var}}
\newcommand{\Cov}{\operatorname{Cov}}
\newcommand{\supp}{\operatorname{supp}}
\DeclareMathOperator*{\argmin}{arg\,min}
\DeclareMathOperator*{\argmax}{arg\,max}
\newcommand{\entrymeta}[3]{{\sffamily\small\color{muted}\textbf{\texttt{#1}}\quad|\quad #2\par #3\par}}
\newcommand{\Problemref}[1]{\texttt{#1}}
\newcommand{\JELref}[2]{\texttt{#2} (JEL \texttt{#1})}
\begin{document}
\section*{Effective bureaucratic organizations}
\textbf{Problem ID:} \texttt{D73-170}\quad \textbf{JEL:} \texttt{D73}\par
% BEGIN ECONOMICS STATEMENT
\label{problem:OP-0382}
\label{audited:ECON-099}
\label{audited:ECON-100}
{\sffamily\small\color{muted}Original subject group: Political economy and institutions\par}
\label{definitions:problem:30007664}
\textbf{Audit classification:} Published research agenda. Full review checked. System-wide reform evaluation is an explicit agenda; the production mapping and finite-menu design are editorial.
\subsection*{Definitions and precise research target}
\textbf{Complete editorial problem.} Fix a public-service system composed of departments \(j=1,\ldots,J\) and officials \(i=1,\ldots,I\). Organizational policy \(p\) from a declared finite menu specifies recruitment, compensation, promotion, delegation, monitoring, and political insulation. Let \(e_{ij}(p)\) be an official's task inputs and \(Q(p)\) measured service quality for the full citizen population. The production mapping \(\widetilde Q(p)=f((e_{ij}(p))_{i,j},p,\eta)\) can contain cross-department complementarities and bottlenecks. Let \(C(p)\) be total cost and \(B\) a fixed administrative budget. Estimate policy effects on \(Q\), distributional access, and corruption, and characterize feasible bundles maximizing a declared quality index subject to \(C(p)\le B\). Testing separate incentive changes is insufficient when one department's performance changes another's feasible output. Identify interactions through system-level reform variation, linked personnel and service data, or a validated organizational model. Address endogenous selection of officials, gaming of measured targets, and changes in political control. State whether any identified optimum applies locally or to the entire system. The cited review calls for studying bureaucracies as interacting systems and evaluating system-wide reforms; the finite policy menu, objective, and interaction estimands here constitute an adapted design problem.

\textbf{Definitions and admissible scope.} Tasks and officials are not interchangeable indices: define assignment set \(\mathcal E\subseteq\{1,\ldots,I\}\times\{1,\ldots,J\}\), effort vector \(e\), and service production \(f(e,p,\eta)\), including resource and departmental constraints. Let citizen outcomes \(q_n(p)\) have declared weights \(\mu_n\), and define \(Q(p)=E[\sum_n\mu_nq_n(p)]\). Corruption and access inequalities enter as separate constraints, not invisible ingredients of a quality score. Recruitment changes who is employed, so a full model specifies entrants and exits and compares citizen outcomes on a fixed population. A known finite policy menu can be optimized by enumeration; the research task is identifying system outcomes, complementarities, target gaming, and transportable organizational counterfactuals. Take finite departments, officials, citizen population, and policy menu. Let the realized population output be \(\widetilde Q(p)=\sum_n\mu_nq_n(p)=f(e(p),p,\eta)\); then the scalar objective is \(Q(p)=E[\widetilde Q(p)]\), resolving the distinction between production realizations and expected quality. Use finite means and explicit budget/access/corruption constraints with stated units. The unknown effort response, shock law, production complementarities, measurement process, and employee selection restrictions are inputs to the admissible model class, which is restricted to match the joint personnel/service data law. Report policy-value and interaction identified sets before presenting conditional maximizers.
\subsection*{Variants and assumption changes}
\begin{enumerate}
\item System-level complementarity (source-motivated): randomize or credibly study linked departmental reforms and estimate \(Q(p_{11})-Q(p_{10})-Q(p_{01})+Q(p_{00})\). Isolated employee incentives cannot identify that contrast.
\item Political control and measurement (source-motivated extension): compare insulation and monitoring bundles across political turnover using independent citizen service measures. Distinguish genuine output from reclassification or gaming of reported targets.
\end{enumerate}
\subsection*{References and source audit}
\textbf{Support and access.} Full review checked. System-wide reform evaluation is an explicit agenda; the production mapping and finite-menu design are editorial. \textbf{Locator:} Section 4; conclusion p. 418, second direction.
\begin{enumerate}\raggedright
\item\hypertarget{definitions:source:30007664:1}{} Timothy Besley; Robin Burgess; Adnan Khan; Guo Xu. 2022. \emph{Bureaucracy and Development}. Section 4, printed pp. 412–413; concluding research agenda, printed p. 418, second direction.
\url{https://www.robinburgess.com/s/Bureaucracy_and_Development.pdf}.
DOI: \url{https://doi.org/10.1146/annurev-economics-080521-011950}.
\end{enumerate}

\subsection*{Integrated refinement: ECON-099}
{\sffamily\small\color{muted}Measuring public-service output rather than bureaucratic inputs\par}
This refinement adopts the additional source conventions
(page~\pageref{audited:conventions}), subject to its local restrictions.
\textbf{Anchored public-service output and incentive-valid measurement.}
For office $j$, let latent quality $Q_j(x)$ follow standardized resource/case-mix
intervention $x=(H,B,X)$ and common reference law $\nu$. With explicit
administrative-record and citizen-validation error/gaming models, validate
$m(Z,X)$ and identify
\[
\theta_{jk}=\int\{\mathbb E[Q_j(x)]-\mathbb E[Q_k(x)]\}\,d\nu(x).
\]
Provide external scale anchors, overlap, sampling and missing-case rules;
arbitrary rescaling changes numerical comparisons. Standardizing inputs or case
mix needs a defensible production/causal model. Record volume, timeliness,
accuracy, coverage and citizen outcomes separately. Validation signals can be
imperfect, and an incentive change may alter recording and gaming. Passive
predictive validity therefore does not establish validity when the measure
determines rewards.
\subsection*{Supplied source audit for ECON-099}
Local [S1], [S2], etc. in this audit and the references immediately below
belong to ECON-099, independently of the original entry's reference numbering.
[S1] and [S2] directly support public-service output measurement and external data. The anchored latent-quality and standardized-office target is editorial; the review does not guarantee recovery of unobserved true quality.
\subsection*{References for integrated refinement ECON-099}
{\footnotesize\raggedright
\par\noindent\textbf{[S1]} Guo Xu, Erika Deserranno, Diana Moreira, and Edoardo Teso (2023). \emph{Bureaucracy}. VoxDevLit 8(1), September 2023.\par \textit{Verified locator / source role:} Primary publisher page checked for review identity, authors, version and background. The specific published agenda is corroborated in [S2]; the exact formal target is editorial.\par \url{https://voxdev.org/voxdevlit/bureaucracy}\par\smallskip
\par\noindent\textbf{[S2]} Oliver Hanney / VoxDev (living compilation). \emph{Evidence gaps: Key unanswered questions in development economics}. Accessed 8 October 2026. The page currently covers 20 topics; the ``16-topics URL is a historical slug.\par \textit{Verified locator / source role:} Topic heading: Bureaucracy. Supports the broad research agenda; this is not a verbatim quotation or an attribution of the displayed mathematical target.\par \url{https://voxdev.org/topic/evidence-gaps-key-unanswered-questions-across-16-topics-development-economics}\par\smallskip
}

\subsection*{Integrated refinement: ECON-100}
{\sffamily\small\color{muted}Adaptive bureaucratic capacity under crises and changing demands\par}
This refinement adopts the additional source conventions
(page~\pageref{audited:conventions}), subject to its local restrictions.
\textbf{Adaptive capacity under unfamiliar shocks.}
Let a nonanticipative policy $\pi\in\Pi$ choose staffing, delegation and
incentives using date-$t$ information. Specify initial resources, budgets,
accountability constraints, normal-service requirements, capacities and
adjustment lags. With quality $Q_t^\pi$, demand $d_t$, real costs $C_t^\pi$ and a
declared monetary benefit conversion, compare
\[
V(\pi)=\mathbb E\sum_{t=0}^{T}\delta^t
              [v_t(Q_t^\pi,d_t)-C_t^\pi].
\]
Identify feasible organizational improvements for stated unfamiliar-shock
scenarios. Preparedness, adaptive staffing and emergency delegation operate at
different dates; expected-value, ambiguity-robust and service-reliability
criteria are distinct. Out-of-support crises need explicit invariance or bounds.
Do not subtract raw quality from dollars or assume within-agency evidence
transports unchanged to another institution; use a supremum unless attainment
is established.
\subsection*{Supplied source audit for ECON-100}
Local [S1], [S2], etc. in this audit and the references immediately below
belong to ECON-100, independently of the original entry's reference numbering.
[S1] and [S2] explicitly support bureaucratic innovation and adaptation to changing/crisis demands. The nonanticipative monetized objective is editorial and makes units, safeguards and transport assumptions explicit.
\subsection*{References for integrated refinement ECON-100}
{\footnotesize\raggedright
\par\noindent\textbf{[S1]} Guo Xu, Erika Deserranno, Diana Moreira, and Edoardo Teso (2023). \emph{Bureaucracy}. VoxDevLit 8(1), September 2023.\par \textit{Verified locator / source role:} Primary publisher page checked for review identity, authors, version and background. The specific published agenda is corroborated in [S2]; the exact formal target is editorial.\par \url{https://voxdev.org/voxdevlit/bureaucracy}\par\smallskip
\par\noindent\textbf{[S2]} Oliver Hanney / VoxDev (living compilation). \emph{Evidence gaps: Key unanswered questions in development economics}. Accessed 8 October 2026. The page currently covers 20 topics; the ``16-topics URL is a historical slug.\par \textit{Verified locator / source role:} Topic heading: Bureaucracy. Supports the broad research agenda; this is not a verbatim quotation or an attribution of the displayed mathematical target.\par \url{https://voxdev.org/topic/evidence-gaps-key-unanswered-questions-across-16-topics-development-economics}\par\smallskip
}

\subsection*{Common conventions: Source mathematical conventions}

These conventions apply to each entry unless explicitly replaced.

\textbf{Models and identification.} A primitive model \(m\) specifies all probability laws, feasible choices, information, equilibrium and measurement rules used by its target. The admissible class \(\mathcal M\) is fixed independently of outcomes. If \(P_O(m)\) is its observable probability law and \(\tau(m)\) the target, its identified set at observable law \(P\) is
\[
 \mathcal I_\tau(P)=\{\tau(m):m\in\mathcal M,\ P_O(m)=P\}.
\]
Point identification means that this set is a singleton. An empty set signals incompatibility of data and assumptions. Coordinate bounds need not describe the joint set. Bounds are sharp only if every retained target is attainable or arbitrarily approachable by compatible models. Finite bounds require explicit restrictions on unobserved quantities; unrestricted confounding, hidden wealth, extrapolation or arbitrary equilibrium selection can yield uninformative sets.

\textbf{Causal interventions.} A potential outcome \(Y_i(p)\) is the outcome under a complete policy regime \(p\), including timing, financing, information and market spillovers. The target population and units are fixed before treatment. With interference, use the complete assignment vector or a justified exposure mapping; individual no-interference notation alone cannot identify an economy-wide intervention. A difference of mean potential outcomes does not require the cross-world joint distribution. Conditioning on post-treatment migration, survival, enrollment or employment changes the estimand and needs separate justification.

\textbf{Statistical statements.} An estimator, confidence set, forecast or test specifies a sampling law, model class and loss/coverage criterion. Uniform \((1-\alpha)\)-coverage means \(\inf_{m\in\mathcal M}\Pr_m\{\tau(m)\in C_n\}\geq1-\alpha\), or an explicitly asymptotic version. The whole parameter space trivially has coverage, so informativeness requires a stated width, risk or power objective. A forecasting improvement is relative to a named benchmark, information set and evaluation population.

\textbf{Equilibrium and optimization.} An equilibrium comprises feasible plans, optimality, beliefs where needed, budget constraints and all market/resource clearing. The primitives or a stated benchmark define each correspondence \(\mathcal E(p;m)\). Nonemptiness is assumed or proved, never inferred just from compact policy sets. A selected-equilibrium target uses a declared selection rule; a robust target quantifies over all permitted equilibria. Use suprema or infima unless attainment is established. An \(\varepsilon\)-optimal policy is feasible and within \(\varepsilon\) of the finite optimum value. Report an empty feasible set separately.

\textbf{Welfare and accounting.} Normative utility, interpersonal normalization, social weights, numeraire, discounting and the use of public receipts are primitives. Behavioral utility need not coincide with the welfare criterion. Household welfare includes ownership income and rents once; adding profits again to compensating variation generally double-counts them. A monetary equivalent is defined by an explicit transfer and price convention, with monotonicity and range conditions. Average effects, mechanism contributions, transfers and welfare losses are different targets. Infinite sums require convergence and dynamic feasible plans require terminal or no-Ponzi conditions.

\textbf{Source versions.} A working-paper date, revision date and journal publication year are distinguished. A DOI for a journal version is not automatically the DOI of an earlier manuscript. Locators refer to the stated version and printed pages where specified. A metadata-only check does not verify a claim about a full-text conclusion. Every entry's source subsection narrows the provenance claim accordingly.

\subsection*{Common conventions: Additional-source status and mathematical conventions}

\label{audited:conventions}

Of the 100 supplied ECON entries, 65 distinct problems appear here; 32 close
matches refine earlier entries, and three duplicate questions map to existing
entries. Appendix~\ref{app:audited-crosswalk} lists every destination.
The attachment's P/R/S/U distinctions and source audits are retained. Its
precise mathematical formalizations are editorial unless the individual audit
says otherwise. Candidate subsidy targets ECON-004--006 retain their U flags;
the resolved approval-core question ECON-007 updates OP-0202 above.
The following supplied conventions govern both this part and the attributed
ECON refinements elsewhere in the catalogue, subject to local restrictions.

\textbf{Parameterized questions.} A problem family lists inputs that must be fixed before an instance is evaluated: population, horizon, policies, information, data law, model class and welfare weights. These are required choices, not quantities the citations are assumed to specify. Undefined applications should not be treated as identified estimands.

\textbf{Indices and integrability.} Sets of agents, firms and regions are finite unless a population measure is explicitly used. \(i,j\) index units, \(r\) regions, \(t\) dates and \(h\) horizons. \(\mathbb E\) and \(\Pr\) refer to the declared population and law; \(\mathbf1\{A\}\) is an indicator. Every displayed expectation and discounted sum needs existence in the stated domain. Infinite horizons use a discount in \((0,1)\) and a convergence condition; finite horizons may use discount one.

\textbf{Optimization and existence.} A displayed \(\max\), \(\min\), or \(\arg\max\) is conditional on attainment. For finite feasible menus this is immediate; general spaces need continuity and compactness/coercivity or another existence argument. Without attainment, the target is the supremum/infimum and arbitrarily near-optimal feasible choices. \(\inf\varnothing=+\infty\). Empty feasibility and an unidentified target are different issues.

\textbf{Potential outcomes.} \(Y_i(z)\) is the outcome under a specified intervention, not the observed conditional mean among people choosing \(z\). In binary comparisons \(\Delta Y_i=Y_i(1)-Y_i(0)\). Specify assignment, treatment versions, compliance, information and observation horizon. Interference is allowed under a full policy vector; compressed exposure notation needs a justified mapping. No interference, consistency and overlap are substantive assumptions, not automatic consequences of notation.

\textbf{Populations and observation.} Fix baseline eligibility and population weights. Conditioning on post-policy employment, survival, relocation or program uptake can change the population being compared. Death is not ordinary loss to follow-up; outcomes truncated by death need a composite convention, joint survival target, or explicitly defined survivor stratum. Fertility-changing interventions need a population functional rather than an unexplained fixed descendant cohort. Measurement and missingness models must be supplied.

\textbf{Identification.} A structural model \(m\in\mathcal M\) implies observable law \(P_m\) and target \(\theta(m)\). Identification means
\[
 P_{m_1}=P_{m_2}\ \Longrightarrow\ \theta(m_1)=\theta(m_2)
 \quad\text{for all }m_1,m_2\in\mathcal M .
\]
Without identification, the target is
\[
 \Theta(P;\mathcal M)=\{\theta(m):m\in\mathcal M,\ P_m=P\}.
\]
Writing \(\theta(P)\) before establishing identification can hide incompatible latent models. Confidence sets additionally need a sampling law, confidence level, asymptotic sequence and uniformity class. Point identification, coverage and informative precision are separate properties.

\textbf{Mediation.} Total effects of randomized assignment are distinct from cross-world natural indirect effects. Belief/effort-channel effects require a meaningful mediator intervention and additional measurement, exclusion or mediation assumptions. Randomization of the initial policy alone does not supply those assumptions. Report bounds or interventional alternatives when the stronger target is unsupported.

\textbf{Equilibrium.} A counterfactual \(W(z)\), \(p(z)\), or \(\mu_t^z\) requires individual optimization, feasibility, government/resource budgets, market clearing and state-transition laws. State equilibrium existence and selection, or report ranges over compatible equilibria. Initial-state integration includes subsequent endogenous transitions; current-state integration must use the policy-dependent distribution. Reduced-form invariance after policy is not assumed.

\textbf{Welfare accounts.} Scores, utility, health, dollars and ecological quantities require explicit conversion before addition. Fees, taxes, subsidies, wages and extortion payments are transfers between parties; distributional weights can make them welfare relevant, but their face value is not automatically a resource loss. State ownership, geographic boundaries, foreign-income weights and fiscal finance. Do not add profits to household welfare already including dividends, or count land capitalization and the underlying benefit twice. A benefit-cost expression must distinguish gross benefits from net welfare and subtract every real cost exactly once.

\textbf{Derivatives and risk.} Log derivatives require positive arguments; local responses require admissible perturbations and specified equilibrium branches. Differentiation under expectations needs regularity/domination, and boundaries may allow only one-sided derivatives. Risk uses a specified law; ambiguity uses a declared nonempty law set on a common history space. Adaptive policies must be nonanticipative. A robust criterion is an editorial choice unless attributed explicitly.

\textbf{Scope overlaps.} Allocation, collective choice, contracts, household bargaining and coordinated policy overlap economics with optimization and game theory. They are retained because they appeared in the original catalogue; no claim of a strictly game-theory-free selection is made.
% END ECONOMICS STATEMENT
\end{document}
